\documentclass[11pt]{article}
\usepackage[T1]{fontenc}
\usepackage{lmodern}
\usepackage{amsmath,amssymb,amsthm}
\usepackage[margin=1in]{geometry}
\usepackage[hidelinks]{hyperref}
\newtheorem{theorem}{Theorem}
\newtheorem{corollary}{Corollary}
\title{A Disproof of Conjecture 00000004607:\\
Convexity, Nonnegativity, and Two Endpoint Zeros}
\author{Independent solution}
\date{October 9, 2026}

\begin{document}
\maketitle

\begin{abstract}
The numerical properties asserted for the concavity correction are mutually
inconsistent on every volume interval $[0,V]$ with $0<V<\infty$. A convex
nonnegative function with zero values at the two endpoints must vanish
throughout the interval, whereas the conjecture requires that only the endpoints
be zeros. This obstruction uses neither symmetry nor a formula for the
correction. The proof is formalized with Mathlib's ordinary real convexity.
\end{abstract}

\section{What the original statement asserts}
The English statement defines the isoperimetric profile as the minimal interface
area needed to enclose a given volume. It then asserts that its concavity
correction is an explicit nonnegative function of volume, ``vanishing only at
zero and at the total volume,'' and that the function is convex and symmetric.
The Chinese statement asserts the same properties; its zero-set clause
specifically refers to zero volume and total volume as the two points where
the function is zero. The parenthetical descriptions in both languages also
refer to two zeros. The complete, unaltered bilingual statement is reproduced
in \texttt{provenance/ORIGINAL.md}.

Consequently, if $V>0$ is the finite total volume and $c$ is the asserted
real-valued function, the following are necessary properties:
\begin{align}
&c\text{ is convex on }[0,V], \label{eq:convex}\\
&c(v)\geq 0\quad(0\leq v\leq V), \label{eq:positive}\\
&c(v)=0\quad\Longleftrightarrow\quad v=0\text{ or }v=V
  \quad(0\leq v\leq V). \label{eq:zeros}
\end{align}
In particular, both endpoint values are zero. The usual volume-complement
symmetry would be $c(V-v)=c(v)$. We include this in the final formal statement,
but the contradiction does not depend on how symmetry is interpreted, because
\eqref{eq:convex}--\eqref{eq:zeros} alone are impossible.

Neither language defines the operation called a concavity correction. We do
not supply a definition of that operation, identify it with a particular
deficit, or assume a relation between it and a chosen profile. Instead, we
disprove necessary properties that the text itself explicitly imposes on its
output. Any proposed definition satisfying the original conclusion would have
to produce a function satisfying \eqref{eq:convex}--\eqref{eq:zeros}.

\section{The obstruction}
\begin{theorem}\label{thm:zero}
Let $V\geq 0$ be a real number. If $c:[0,V]\to\mathbb R$ is convex and
nonnegative and $c(0)=c(V)=0$, then $c(v)=0$ for every $v\in[0,V]$.
\end{theorem}
\begin{proof}
If $V=0$, the interval has only its zero endpoint. If $V>0$, fix $v\in[0,V]$
and set $t=v/V\in[0,1]$. Then $v=(1-t)0+tV$, so convexity gives
\[
 c(v)\leq(1-t)c(0)+tc(V)=0.
\]
Nonnegativity gives the reverse inequality, hence $c(v)=0$.
\end{proof}

\begin{corollary}\label{cor:no}
For every real $V>0$, no real-valued function on $[0,V]$ satisfies
\eqref{eq:convex}--\eqref{eq:zeros}. In particular, no explicit symmetric
function satisfies all of the properties in the conjecture.
\end{corollary}
\begin{proof}
Theorem~\ref{thm:zero} implies $c(V/2)=0$. But $0<V/2<V$, contradicting
\eqref{eq:zeros}. Requiring a formula or symmetry cannot restore existence.
\end{proof}

Equivalently, just the midpoint case of convexity already gives
\[
 0\leq c(V/2)\leq \tfrac12c(0)+\tfrac12c(V)=0,
\]
while the prescribed zero set and nonnegativity require $c(V/2)>0$.
There are no smoothness, continuity, differentiability, or strict-convexity
assumptions in this argument.

\section{Why the volume interval is a legitimate setting}
Finite positive total volume is the ordinary nondegenerate interpretation of
the text's two endpoint volumes. It is not a curvature hypothesis or a
geometric restriction introduced to obtain a special formula. For example,
the flat square torus $(\mathbb R/\mathbb Z)^2$, with its standard product
area measure, has total area $1$. For every $v\in[0,1]$, the image of
$[0,v)\times[0,1)$ has area $v$. For $0<v<1$ it is a strip bounded by two
circles of finite total length $2$; at the endpoints one uses the empty set
and the whole torus. Thus every volume in $[0,1]$ is admissible in a standard
compact, connected, smooth isoperimetric setting. In the usual infimum
formulation of the profile, these competitors make its value finite and
nonnegative. No computation of the profile or existence theorem for
minimizers is used in the disproof.

More abstractly, if the domain of volumes is convex and contains $0$ and $V$,
it necessarily contains the whole interval between them, in particular $V/2$.
This is also built into Mathlib's \texttt{ConvexOn} predicate. The claimed
convexity therefore cannot evade the midpoint merely by discarding interior
volumes. Restricting the domain to the two endpoints would change the usual
meaning of convexity and would not describe the example above.

When $V=0$, the two endpoints coincide and the identically zero function on
the singleton interval presents no contradiction. This degenerate case cannot
establish a claim about two distinct zeros in positive volume. Infinite total
volume is not a second real endpoint, and the formula $V-v$ is then not a
real-valued reflection of a bounded interval. The original text supplies no
alternative definitions for that case. Our disproof concerns every finite
positive total volume, which includes the concrete ordinary setting just
described. It makes no claim that all conceivable isoperimetric spaces have
finite volume. In particular, no infinite-volume extension can rescue the
asserted conclusion in this finite-volume setting.

\clearpage
\section{Lean formalization and its boundary}
The project is pinned to Lean 4.19.0 and Mathlib commit
\begin{center}
\small\texttt{c44e0c8ee63ca166450922a373c7409c5d26b00b}.
\end{center}
\texttt{Disproof.lean} contains these theorems in the namespace
\texttt{IsoperimetricCorrection}:
\begin{itemize}
\item \texttt{zero\_on\_volume\_interval}: Theorem~\ref{thm:zero}, for every
real $V\geq0$.
\item \texttt{no\_nonnegative\_convex\_double\_zero}:\\
The nonexistence in
Corollary~\ref{cor:no}, for every real~\mbox{$V>0$}.
\item \texttt{no\_shape\_requirements}: nonexistence with the additional
condition $c(V-v)=c(v)$, again for every real~\mbox{$V>0$}.
\end{itemize}
The definition \texttt{ShapeRequirements} groups these explicitly stated
numerical properties; it is not a definition of a concavity correction.
The function is represented as $f:\mathbb R\to\mathbb R$, with every
hypothesis restricted to $[0,V]$. Values outside the interval are immaterial;
any function defined only on the interval can be extended there arbitrarily.

The formal proof uses the stock predicates \texttt{ConvexOn} and
\texttt{Set.Icc}. It obtains the upper bound on each value from
\texttt{ConvexOn.le\_on\_segment} and \texttt{Icc\_subset\_segment}, combines
it with nonnegativity, and applies the exact zero-set condition at $V/2$.
The only positivity hypothesis on the total volume in the final theorems is
$0<V$. No hypothesis asserts the desired conclusion or a geometric theorem
that has not been established.

The Lean project does not formalize the geometry of the torus, the
isoperimetric profile, a specific correction operator, or an interpretation
of ``explicit.'' These are unnecessary for the universal scalar obstruction;
the torus paragraph only verifies in ordinary mathematics that the volume
domain has a familiar geometric realization. The logical passage from the
bilingual prose to the necessary scalar properties is explained above,
rather than hidden as an extra axiom.

\section{Verification and provenance}
The entire project was built against pre-existing stock Mathlib dependencies;
this is not a claim that all of Mathlib was rebuilt. Each source file was also
replayed with \texttt{warningAsError=true}. The axiom audit of all three final
theorems reports only \texttt{propext}, \texttt{Classical.choice}, and
\texttt{Quot.sound}. There are no admitted proofs, custom axioms, or uses of
\texttt{native\_decide}. Build outputs, exit codes, and the unsuccessful
development attempts are preserved under \texttt{verification/} and
\texttt{provenance/}, as described in \texttt{README.md}.

The author input was restricted to the original bilingual statement, both
competition rule files, and stock Lean configuration. The statement's SHA-256
is recorded in \texttt{provenance/INPUTS.json}. No prior solution, public claim
history, selector material, or other author project was used as mathematical
input. The mathematical argument is the elementary endpoint inequality above;
there are no auxiliary numerical computations or external mathematical sources.
\end{document}
